\subsection{黄金单位障碍与精确谱量的超越性}

正文已经分别建立了整数阶 Rényi 维数的“有理或超越”二分、\(q=4\) 与
\(q=5\) 的五次 Perron 数域结构、以及共振窗 \(q=16,17\) 上最小递推特征多项式与 Newman 比值的范数律。
把这些结果在黄金单位 \(\varphi=(1+\sqrt5)/2\) 的有理幂约束下并读后，可进一步抽出一条算术后果链：
虽然 Fold 的高阶压力斜率在极限上由 \(\varphi\) 锁定，但若把精确谱值试图压到 \(\varphi\) 的代数幂轨道上，则会立即遭遇代数单位障碍。

\begin{theorem}[黄金单位的有理幂障碍]\label{thm:derived-fold-golden-rational-power-unit-obstruction}
令
\[
\varphi:=\frac{1+\sqrt5}{2}.
\]
若 \(t\in\QQ\)，则 \(\varphi^t\) 必为代数单位。
特别地，若某个代数数 \(u\) 在其生成数域 \(K=\QQ(u)\) 中满足
\[
\mathrm N_{K/\QQ}(u)\neq \pm 1,
\]
则 \(u\) 不可能写成 \(\varphi^t\)（其中 \(t\in\QQ\)）。
\end{theorem}
\leanverified{paper\_derived\_fold\_golden\_rational\_power\_unit\_obstruction}

\begin{proof}
把 \(t=a/b\in\QQ\) 写成既约分数，其中 \(b\ge 1\)。
先设 \(a\ge 0\)，并记 \(\psi=(1-\sqrt5)/2=-\varphi^{-1}\)。
令 \(u=\varphi^{a/b}\)，则 \(u^b=\varphi^a\)，从而
\[
(u^b-\varphi^a)(u^b-\psi^a)
=
u^{2b}-(\varphi^a+\psi^a)u^b+(\varphi\psi)^a.
\]
由于 \(\varphi^a+\psi^a=L_a\in\ZZ\) 为 Lucas 数，且 \(\varphi\psi=-1\)，故 \(u\) 满足首一整系数方程
\[
X^{2b}-L_aX^b+(-1)^a=0.
\]
因此 \(u\) 为代数整数，且其最小多项式常数项整除 \((-1)^a\)，故必为 \(\pm1\)。
于是 \(u\) 为代数单位。

若 \(a<0\)，则
\[
\varphi^{a/b}=\bigl(\varphi^{-a/b}\bigr)^{-1},
\]
而上一步已知 \(\varphi^{-a/b}\) 为代数单位，故其逆元仍为代数单位。
因此对任意 \(t\in\QQ\) 都有 \(\varphi^t\) 为代数单位。最后一句只是其逆否命题。证毕。
\end{proof}

\leanverified{paper\_derived\_fold\_renyi\_dimension\_transcendence\_q4\_q5\_q16\_q17}
\begin{theorem}[显式阶上的 Rényi 维数超越性]\label{thm:derived-fold-renyi-dimension-transcendence-q4-q5-q16-q17}
沿用命题 \ref{prop:pom-renyi-dimension-spectrum} 的记号，记
\[
D_q=\frac{q\log 2-\log r_q}{(q-1)\log\varphi}.
\]
则
\[
D_4,\qquad D_5,\qquad D_{16},\qquad D_{17}
\]
全部为超越数。
\end{theorem}

\begin{proof}
由定理 \ref{thm:pom-renyi-dimension-no-algebraic-irrationality}，只需排除这四个 \(D_q\) 落在 \(\QQ\) 中。
反设 \(D_q\in\QQ\)，并令
\[
u_q:=\frac{2^q}{r_q}.
\]
由定义
\[
u_q=\exp\!\bigl((q-1)D_q\log\varphi\bigr)=\varphi^{(q-1)D_q}.
\]
由于 \((q-1)D_q\in\QQ\)，定理
\ref{thm:derived-fold-golden-rational-power-unit-obstruction} 蕴含 \(u_q\) 必为代数单位。

下面对四个 \(q\) 分别计算 \(\mathrm N(u_q)\)。
对 \(q=4\)，由推论 \ref{cor:pom-s4-asymptotic} 与命题 \ref{prop:pom-a4-galois-s5}，
\[
r_4
\]
是五次不可约多项式
\[
P_4(\lambda)=\lambda^5-2\lambda^4-7\lambda^3-2\lambda+2
\]
的 Perron 根，故
\[
\mathrm N_{\QQ(r_4)/\QQ}(r_4)=(-1)^5P_4(0)=-2.
\]
于是
\[
\mathrm N(u_4)=\frac{2^{20}}{-2}=-2^{19}\neq \pm1.
\]

对 \(q=5\)，由命题 \ref{prop:pom-s5-recurrence} 与命题 \ref{prop:pom-s5-galois-s5}，
\[
r_5
\]
是五次不可约多项式
\[
P_5(\lambda)=\lambda^5-2\lambda^4-11\lambda^3-8\lambda^2-20\lambda+10
\]
的 Perron 根，故
\[
\mathrm N_{\QQ(r_5)/\QQ}(r_5)=(-1)^5P_5(0)=-10.
\]
于是
\[
\mathrm N(u_5)=\frac{2^{25}}{-10}=-\frac{2^{24}}{5}\neq \pm1.
\]

对 \(q=16,17\)，由命题 \ref{prop:pom-resonance-window-minpoly-irreducible-q9-17}
与注 \ref{rem:pom-moment-charpoly-16-17}，\(r_{16},r_{17}\) 的最小多项式分别为
\[
P_{16}(\lambda),\qquad P_{17}(\lambda),
\]
次数都等于 \(13\)，且常数项满足
\[
P_{16}(0)=8,\qquad P_{17}(0)=673948170.
\]
故
\[
\mathrm N_{\QQ(r_{16})/\QQ}(r_{16})=(-1)^{13}P_{16}(0)=-8,
\qquad
\mathrm N_{\QQ(r_{17})/\QQ}(r_{17})=(-1)^{13}P_{17}(0)=-673948170.
\]
从而
\[
\mathrm N(u_{16})=\frac{2^{208}}{-8}=-2^{205}\neq \pm1,
\]
\[
\mathrm N(u_{17})=\frac{2^{221}}{-673948170}\neq \pm1.
\]

这与 \(u_q\) 为代数单位矛盾。
故 \(D_4,D_5,D_{16},D_{17}\notin\QQ\)。
再由定理 \ref{thm:pom-renyi-dimension-no-algebraic-irrationality} 的二分律，它们只能是超越数。证毕。
\end{proof}

\begin{corollary}[Perron 黄金指数的超越性]\label{cor:derived-fold-perron-golden-exponent-transcendence-q4-q5-q16-q17}
对 \(q\in\{4,5,16,17\}\)，定义
\[
\Theta_q:=\frac{\log r_q}{\log\varphi}.
\]
则
\[
\Theta_4,\qquad \Theta_5,\qquad \Theta_{16},\qquad \Theta_{17}
\]
全部为超越数。
\end{corollary}

\begin{proof}
若 \(\Theta_q\in\QQ\)，则 \(r_q=\varphi^{\Theta_q}\) 为 \(\varphi\) 的有理幂。
由定理 \ref{thm:derived-fold-golden-rational-power-unit-obstruction}，\(r_q\) 必为代数单位。
但在定理 \ref{thm:derived-fold-renyi-dimension-transcendence-q4-q5-q16-q17} 的证明中已算得
\[
\mathrm N_{\QQ(r_4)/\QQ}(r_4)=-2,\qquad
\mathrm N_{\QQ(r_5)/\QQ}(r_5)=-10,
\]
\[
\mathrm N_{\QQ(r_{16})/\QQ}(r_{16})=-8,\qquad
\mathrm N_{\QQ(r_{17})/\QQ}(r_{17})=-673948170,
\]
均不等于 \(\pm1\)，矛盾。

若 \(\Theta_q\) 是代数但非有理，则由 Gelfond--Schneider 定理，\(\varphi^{\Theta_q}=r_q\) 必为超越数；
而 \(r_q\) 已知是代数数，再度矛盾。
故 \(\Theta_q\) 只能是超越数。证毕。
\end{proof}

\begin{theorem}[Newman 比值黄金指数的超越性]\label{thm:derived-fold-newman-golden-exponent-transcendence-q4-q16-q17}
在已核验的 \(q\in\{4,16,17\}\) 上，沿用正文的 Newman 比值记号
\[
\upsilon_q:=\frac{\Lambda_q^2}{r_q},
\qquad
\Xi_q:=\frac{\log \upsilon_q}{\log\varphi}.
\]
则
\[
\Xi_4,\qquad \Xi_{16},\qquad \Xi_{17}
\]
全部为超越数。
\end{theorem}

\begin{proof}
先排除 \(\Xi_q\in\QQ\)。
若 \(\Xi_q\in\QQ\)，则 \(\upsilon_q=\varphi^{\Xi_q}\) 为 \(\varphi\) 的有理幂。
由定理 \ref{thm:derived-fold-golden-rational-power-unit-obstruction}，\(\upsilon_q\) 必为代数单位。

但命题 \ref{prop:pom-a4-newman-u4-algebraic-certificate} 给出
\[
\mathrm N_{\QQ(\upsilon_4)/\QQ}(\upsilon_4)=16,
\]
而定理 \ref{thm:pom-resonance-newman-uq-degree-norm-q16-q17} 给出
\[
\mathrm N_{\QQ(\upsilon_{16})/\QQ}(\upsilon_{16})=8^{12},
\qquad
\mathrm N_{\QQ(\upsilon_{17})/\QQ}(\upsilon_{17})=673948170^{12}.
\]
三者都不等于 \(\pm1\)，故 \(\Xi_q\notin\QQ\)。

若 \(\Xi_q\) 是代数但非有理，则由 Gelfond--Schneider 定理，
\[
\varphi^{\Xi_q}=\upsilon_q
\]
必为超越数；然而 \(\upsilon_q\) 已知为代数数，矛盾。
故 \(\Xi_4,\Xi_{16},\Xi_{17}\) 必为超越数。证毕。
\end{proof}

\begin{corollary}[精确谱量与黄金幂轨道的代数失配]\label{cor:derived-fold-no-algebraic-golden-power-locking}
在上述已核验的阶数上，以下三类精确谱量都不落在 \(\varphi\) 的任何代数幂轨道上：
\[
r_q\neq \varphi^\alpha
\qquad
(\alpha\in\overline{\QQ},\ q\in\{4,5,16,17\}),
\]
\[
\frac{2^q}{r_q}\neq \varphi^\beta
\qquad
(\beta\in\overline{\QQ},\ q\in\{4,5,16,17\}),
\]
\[
\upsilon_q\neq \varphi^\gamma
\qquad
(\gamma\in\overline{\QQ},\ q\in\{4,16,17\}).
\]
换言之，虽然这些精确谱量在极限斜率上受 \(\varphi\) 控制，但在代数层面却系统性地拒绝锁相到黄金单位所生成的代数幂轨道。
\end{corollary}

\begin{proof}
对 \(r_q\) 的断言正是推论
\ref{cor:derived-fold-perron-golden-exponent-transcendence-q4-q5-q16-q17}
的改写。
对 \(2^q/r_q\) 而言，若存在 \(\beta\in\overline{\QQ}\) 使
\(
2^q/r_q=\varphi^\beta
\)，则当 \(\beta\in\QQ\) 时与定理
\ref{thm:derived-fold-renyi-dimension-transcendence-q4-q5-q16-q17}
证明中的非单位障碍矛盾；当 \(\beta\notin\QQ\) 但仍代数时，又与 Gelfond--Schneider 定理矛盾，因为左端是代数数。
对 \(\upsilon_q\) 的情形同理，只需改用定理
\ref{thm:derived-fold-newman-golden-exponent-transcendence-q4-q16-q17}
中的非单位障碍与 Gelfond--Schneider 定理。证毕。
\end{proof}

\begin{conjecture}[全阶非单位刚性与全谱超越性]\label{conj:derived-fold-nonunit-rigidity-full-spectrum}
对每个整数 \(q\ge 2\)，设 \(r_q\) 为 Fold moment-kernel 的精确 Perron 根，并定义
\[
u_q:=\frac{2^q}{r_q}.
\]
当 Newman 比值 \(\upsilon_q=\Lambda_q^2/r_q\) 有定义时，进一步考虑 \(\upsilon_q\)。
则预期始终有
\[
u_q\ \text{不是代数单位},
\qquad
\upsilon_q\ \text{不是代数单位}.
\]
在该断言下，定理 \ref{thm:pom-renyi-dimension-no-algebraic-irrationality} 与定理
\ref{thm:derived-fold-golden-rational-power-unit-obstruction}
立即推出：对所有整数 \(q\ge 2\)，Rényi 维数
\[
D_q=\frac{q\log 2-\log r_q}{(q-1)\log\varphi}
\]
恒为超越数；而在 \(\upsilon_q\) 有定义的全部阶数上，
\[
\frac{\log \upsilon_q}{\log\varphi}
\]
亦恒为超越数。
\end{conjecture}

\begin{remark}[黄金斜率与非 Kummer 精确谱值]\label{rem:derived-fold-golden-slope-vs-nonkummer-exactness}
这一组推论把正文中“\(\frac1q\log r_q\to \frac12\log\varphi\)”的黄金斜率，与精确谱值的代数性质严格分离开来：
极限热力学斜率可以是黄金的，而精确 Perron 根、归一化量 \(2^q/r_q\) 以及 Newman 比值 \(\upsilon_q\) 却仍然不落在任何 \(\varphi\) 的代数幂轨道上。
因此，Fold 精确谱的算术指纹并不是“黄金单位上的代数锁相”，而更接近一种稳定的非单位刚性。
\end{remark}

\endinput
